% ==STYLE-BEGIN==
\documentclass[UTF8,11pt]{ctexart}
\usepackage{amsmath,amssymb,amsthm}
\usepackage[a4paper,margin=2.4cm,top=2.3cm,bottom=2.3cm]{geometry}
\usepackage{fancyhdr}
\usepackage{enumitem}
\usepackage{hyperref}
\hypersetup{hidelinks}

% ---- 版式 ----
\linespread{1.15}
\setlength{\parskip}{0.35em}
\setlength{\parindent}{1.6em}
\setlist{nosep,leftmargin=2em}

% ---- 页眉 ----
\pagestyle{fancy}
\fancyhf{}
\fancyhead[L]{\small\coursename}
\fancyhead[R]{\small\thepage}
\renewcommand{\headrulewidth}{0.4pt}

% ---- 定理环境：中文用正体，标题加粗，按小节共用一条流水号 ----
\newtheoremstyle{cnthm}{0.7em}{0.7em}{\normalfont}{}{\bfseries}{}{0.5em}{}
\theoremstyle{cnthm}
\newtheorem{thm}{定理}[section]
\newtheorem{lem}[thm]{引理}
\newtheorem{prop}[thm]{命题}
\newtheorem{cor}[thm]{推论}
\newtheorem{defn}[thm]{定义}
\newtheorem{rem}[thm]{注}
\newtheorem{exa}[thm]{例}
\newtheorem{conj}[thm]{猜想}
\renewcommand{\proofname}{\bfseries 证明}
\providecommand{\coursename}{}
% ==STYLE-END==
\renewcommand{\coursename}{高等代数\,·\,1118}
\xeCJKDeclareCharClass{CJK}{"2150 -> "218F}
\title{高等代数 1118}
\date{}
\begin{document}
\maketitle

\section{矩阵的运算: 乘法}

设 $A=(a_{ij})\in M_{m,n}(K)$, $B=(b_{ij})\in M_{n,s}(K)$, 则 $AB\in M_{m,s}(K)$, 且
\begin{equation}
(AB)(i,j)=(a_{i1},\cdots,a_{in})
\begin{pmatrix}b_{1j}\\\vdots\\b_{nj}\end{pmatrix}
=\sum_{k=1}^{n}a_{ik}b_{kj}.
\tag{1}
\end{equation}

\begin{thm}[乘法的性质]
(1) $(AB)C=A(BC)$(结合律);

(2) $A(B+C)=AB+AC$(分配律);

(3) $k(AB)=(kA)B=A(kB)$;

(4) $(AB)^{\mathrm T}=B^{\mathrm T}A^{\mathrm T}$, $(A+B)^{\mathrm T}=A^{\mathrm T}+B^{\mathrm T}$.
\end{thm}

\begin{rem}[三条"没有"]
(1) $AB\neq BA$: \textbf{无交换律};

(2) $AB=O$ 推不出 $A=O$ 或 $B=O$: 有\textbf{零因子};

(3) $AB=AC$ 且 $A\neq O$ 推不出 $B=C$: \textbf{无消去律}.
\end{rem}

\section{矩阵乘法的几何意义}

设 $A=(a_{ij})\in M_{m,n}(K)$, 定义映射
\[
f_A:\ K^n\longrightarrow K^m,\qquad X\longmapsto AX .
\]

\textbf{$f_A$ 保持运算}:
\[
f_A(X+Y)=A(X+Y)=AX+AY=f_A(X)+f_A(Y),
\]
\[
f_A(kX)=A(kX)=k(AX)=k\,f_A(X).
\]
于是 $f_A$ 是一个\textbf{线性映射}.

记 $\varepsilon_1,\cdots,\varepsilon_n$ 是 $K^n$ 的标准基, 则有
\[
f_A(\varepsilon_j)=A\varepsilon_j=A_j,
\]
其中 $A_j$ 是 $A$ 的第 $j$ 个列向量, 即: \textbf{第 $j$ 列就是第 $j$ 个基向量的像}.

\textbf{反过来}, 若 $g:K^n\to K^m$ 是线性映射, 记 $B_j=g(\varepsilon_j)$, 则 $\forall X=(x_1,\cdots,x_n)^{\mathrm T}$ 有
\begin{equation}
g(X)=g\Bigl(\sum_{j=1}^{n}x_j\varepsilon_j\Bigr)
=\sum_{j=1}^{n}x_j\,g(\varepsilon_j)
=\sum_{j=1}^{n}x_jB_j
=(B_1,B_2,\cdots,B_n)X .
\tag{2}
\end{equation}
这说明: \textbf{$K^n$ 到 $K^m$ 的线性映射与 $M_{m,n}(K)$ 中的矩阵一一对应}.

\begin{exa}[用映射看方程组]
线性方程组 $AX=\beta$ 就是
\[
f_A(X)=\beta,
\]
即"求 $X$, 使它在 $f_A$ 下的像等于 $\beta$". 而
\[
f_Af_B(X)=f_A(BX)=A(BX)=(AB)X,
\]
"先作用 $B$ 再作用 $A$"对应矩阵乘积 $AB$.

于是两个线性映射的复合的矩阵等于它们对应矩阵的乘积:
\[
\sigma:\ L(K^n,K^m)\longrightarrow M_{m,n}(K),
\]
$\sigma$ 是双射, 且 $\sigma(fg)=\sigma(f)\sigma(g)$.
\end{exa}

\subsection{从映射复合看矩阵乘法结合律}

映射的复合天然满足结合律:
\[
f_A\circ(f_B\circ f_C)=(f_A\circ f_B)\circ f_C .
\]
用 $\sigma$ 翻译成矩阵:
\[
A(BC)=\sigma\bigl(f_A\circ(f_B\circ f_C)\bigr)
=\sigma\bigl((f_A\circ f_B)\circ f_C\bigr)=(AB)C .
\]
所以\textbf{矩阵乘法的结合律不是"算"出来的}, 而是"映射复合"这个操作本身的性质——这也解释了它为什么"看起来显然却算起来麻烦".

\subsection{例: 旋转矩阵}

设 $R_\theta:\mathbb{R}^2\to\mathbb{R}^2$ 把向量 $X$ 逆时针旋转 $\theta$ 角, $X=(x,y)^{\mathrm T}$. 由"第 $j$ 列 $=$ 第 $j$ 个基向量的像",
\[
R_\theta(\varepsilon_1)=\cos\theta\,\varepsilon_1+\sin\theta\,\varepsilon_2,
\qquad
R_\theta(\varepsilon_2)=-\sin\theta\,\varepsilon_1+\cos\theta\,\varepsilon_2,
\]
故
\[
R_\theta(X)=x\,R_\theta(\varepsilon_1)+y\,R_\theta(\varepsilon_2)
=\begin{pmatrix}\cos\theta&-\sin\theta\\ \sin\theta&\cos\theta\end{pmatrix}
\begin{pmatrix}x\\y\end{pmatrix}.
\]
把两个旋转复合:
\[
\begin{pmatrix}\cos\theta_2&-\sin\theta_2\\ \sin\theta_2&\cos\theta_2\end{pmatrix}
\begin{pmatrix}\cos\theta_1&-\sin\theta_1\\ \sin\theta_1&\cos\theta_1\end{pmatrix}
=
\begin{pmatrix}\cos(\theta_1+\theta_2)&-\sin(\theta_1+\theta_2)\\
\sin(\theta_1+\theta_2)&\cos(\theta_1+\theta_2)\end{pmatrix},
\]
即 $R_{\theta_2}R_{\theta_1}=R_{\theta_1+\theta_2}$: 几何上"转两次等于转一次", 代数上就是这两个矩阵相乘.

\section{矩阵的秩不等式}

\begin{thm}
若向量组 $\alpha_1,\cdots,\alpha_s$ 能被向量组 $\beta_1,\cdots,\beta_t$ 线性表示, 则
\[
\operatorname{rank}(\alpha_1,\cdots,\alpha_s)\leq\operatorname{rank}(\beta_1,\cdots,\beta_t).
\]
\end{thm}

\begin{proof}
用子空间的语言: "能被表示"就是说
\[
L(\alpha_1,\cdots,\alpha_s)\subseteq L(\beta_1,\cdots,\beta_t),
\]
而子空间的维数不超过包含它的子空间的维数, 故
\[
\dim L(\alpha_1,\cdots,\alpha_s)\leq\dim L(\beta_1,\cdots,\beta_t),
\]
即 $\operatorname{rank}(\alpha_1,\cdots,\alpha_s)\leq\operatorname{rank}(\beta_1,\cdots,\beta_t)$.

也可以只用向量组语言: 取 (Ⅱ) 的极大无关组 $(\text{Ⅱ}')$, 则 (Ⅰ) 能被 $(\text{Ⅱ}')$ 线性表示; 于是 (Ⅰ) 的极大无关组能被 $(\text{Ⅱ}')$ 表示, 而它线性无关, 故所含向量个数不超过 $|\text{Ⅱ}'|$.
\end{proof}

\textbf{$AB$ 的列向量是 $A$ 的列向量的线性组合}:
\[
AB=(A_1,\cdots,A_n)
\begin{pmatrix}b_{11}&b_{12}&\cdots&b_{1s}\\
b_{21}&b_{22}&\cdots&b_{2s}\\
\vdots&\vdots&&\vdots\\
b_{n1}&b_{n2}&\cdots&b_{ns}\end{pmatrix}
=\Bigl(\sum_{k=1}^{n}b_{k1}A_k,\ \cdots,\ \sum_{k=1}^{n}b_{ks}A_k\Bigr).
\]
于是 $AB$ 的列组能被 $A$ 的列组线性表示, 由命题 1 得 $r(AB)\leq r(A)$; 又
\[
r(AB)=r\bigl((AB)^{\mathrm T}\bigr)=r(B^{\mathrm T}A^{\mathrm T})\leq r(B^{\mathrm T})=r(B),
\]
合起来:

\begin{thm}
$r(AB)\leq\min\{r(A),\ r(B)\}$.
\end{thm}

\begin{thm}
若 $k\neq0$, 则 $r(kA)=r(A)$.
\end{thm}

\begin{proof}
$kA$ 的列组是 $kA_1,\cdots,kA_n$, 与 $A_1,\cdots,A_n$ 线性等价, 故秩相同.
\end{proof}

\begin{thm}
$r(A+B)\leq r(A)+r(B)$.
\end{thm}

\begin{proof}
设 $A=(A_1,\cdots,A_n)$, $B=(B_1,\cdots,B_n)$, 取 $A$ 的列组的极大无关组 $A_{i_1},\cdots,A_{i_r}$, $B$ 的列组的极大无关组 $B_{j_1},\cdots,B_{j_t}$. 则 $A+B$ 的每一列 $A_k+B_k$ 都能被这 $r+t$ 个向量线性表示, 故
\[
r(A+B)\leq r+t=r(A)+r(B).
\]
\end{proof}

\begin{thm}[Sylvester 秩不等式]
设 $A\in M_{m,n}(K)$, $B\in M_{n,s}(K)$, 则
\[
r(AB)\geq r(A)+r(B)-n .
\]
\end{thm}

\begin{proof}
记 $AB=C$, $B=(B_1,\cdots,B_s)$, $C=(C_1,\cdots,C_s)$, 则
\[
A(B_1,\cdots,B_s)=(C_1,\cdots,C_s)\quad\Longrightarrow\quad AB_j=C_j\quad(j=1,\cdots,s),
\]
即每个 $B_j$ 都是方程组 $AX=C_j$ 的解.

取 $C_1,\cdots,C_s$ 的极大无关组 $C_{i_1},\cdots,C_{i_r}$($r=r(C)$), 则每个 $C_j$ 可写成
\[
C_j=k_1C_{i_1}+\cdots+k_rC_{i_r},
\]
于是
\[
A\bigl(k_1B_{i_1}+\cdots+k_rB_{i_r}\bigr)=k_1C_{i_1}+\cdots+k_rC_{i_r}=C_j .
\]
这说明
\[
B_j-\bigl(k_1B_{i_1}+\cdots+k_rB_{i_r}\bigr)
\]
是齐次方程组 $AX=O$ 的解. 取 $AX=O$ 的基础解系 $\eta_1,\cdots,\eta_{n-r(A)}$, 则每个 $B_j$ 都能被
\[
B_{i_1},\cdots,B_{i_r},\ \eta_1,\cdots,\eta_{n-r(A)}
\]
这 $r+n-r(A)$ 个向量线性表示, 故
\[
r(B)\leq r+n-r(A)=r(AB)+n-r(A),
\]
移项即得 $r(AB)\geq r(A)+r(B)-n$.
\end{proof}

\end{document}
